\section{Monotonicity-based guard erasure}
\label{sec:erasure}

Guards are the tax the untyped encoding pays for soundness.  This
section determines exactly when the tax can be waived: for which
types $\tau$ the guard atoms $\guard{\tau}(\cdot)$ may be deleted
from a fully monomorphised, fully specialised problem without
compromising the soundness invariant.  The answer --- $\tau$ must be
\emph{witnessed inhabited} and \emph{monotonic}, with monotonicity
established either by the absence of naked variables or by a
provable-infinity criterion that the translation's own injectivity
and discrimination axioms supply for recursive inductive types ---
transplants the monotonicity framework of Claessen, Lillieström and
Smallbone~\cite{Claessen2011} and Blanchette, B\"ohme, Popescu and
Smallbone~\cite{Blanchette2016} from many-sorted first-order logic
to the guarded unsorted problems generated from $\CICm$.  The
transplant is not routine: in the guarded world all types share one
domain, an individual may inhabit several types, and the empty type
is expressible.  Both phenomena show up as necessary side conditions
(\cref{prop:nec-inhab,prop:nec-mono}).

Throughout this section, formulas are in negation normal form
(classically without loss of generality), and $(\Phi, \gamma)$ is a
totally specialised problem (\cref{def:spec,lem:spec-total}); we
write $\mathcal{G} = \setof{\guard{\tau}}{\tau \in \mathcal{T}}$ for
its finite family of guard predicates and $R \coloneqq \Phi \cup
\{\mathrm{nnf}(\neg\gamma)\}$ for the refutation set.

\subsection{Erasure, and what must be proved}

\begin{definition}[Erasure]
\label{def:erasure}
For $\Sigma_0 \subseteq \mathcal{T}$, the erasure
$\erase{\Sigma_0}(\varphi)$ replaces every atom $\guard{\tau}(t)$
with $\tau \in \Sigma_0$ by $\top$ (in NNF: positive occurrences by
$\top$, negative occurrences $\neg\guard{\tau}(t)$ by $\bot$),
followed by the evident Boolean simplifications.  In particular
guarded quantifiers lose their guards: $\forall x.\,
(\neg\guard{\tau}(x) \vee \varphi) \mapsto \forall x.\, \varphi$ and
$\exists x.\, (\guard{\tau}(x) \wedge \varphi) \mapsto \exists x.\,
\varphi$, and typing-axiom conclusions $\guard{\tau}(f(\tup{x}))$
become trivial.  $\erase{\Sigma_0}$ is applied to every element of
$R$.
\end{definition}

\begin{lemma}[Erasure at total guards]
\label{lem:erasure-total}
If $\FM$ is a structure in which $\guard{\tau}^{\FM}$ is the whole
domain for every $\tau \in \Sigma_0$, then $\FM \folmodels \varphi$
iff $\FM \folmodels \erase{\Sigma_0}(\varphi)$, for every formula
$\varphi$.
\end{lemma}

\begin{proof}
Substituting $\top$ for atoms that are true under every valuation
does not change the truth value of any formula; by induction on
$\varphi$.
\end{proof}

\begin{proposition}[Erasure never loses proofs]
\label{prop:erasure-complete}
If $\Phi \folproves \gamma$ then $\erase{\Sigma_0}(\Phi) \folproves
\erase{\Sigma_0}(\gamma)$, for any $\Sigma_0$.
\end{proposition}

\begin{proof}
Contrapositively, let $\FM \folmodels \erase{\Sigma_0}(R)$.  Let
$\FM'$ agree with $\FM$ except that $\guard{\tau}^{\FM'}$ is the
whole domain for each $\tau \in \Sigma_0$.  The erased formulas do
not mention the reinterpreted predicates, so $\FM' \folmodels
\erase{\Sigma_0}(R)$, and by \cref{lem:erasure-total} $\FM'
\folmodels R$.
\end{proof}

So erasure is dangerous in one direction only: it may \emph{prove
more}.  The soundness obligation is:

\begin{definition}[Inflatability]
\label{def:inflatable}
$R$ is \emph{$\Sigma_0$-inflatable} if, whenever $R$ is
satisfiable, $R$ has a model in which $\guard{\tau}$ is total (the
whole domain) for every $\tau \in \Sigma_0$.
\end{definition}

\begin{theorem}[Erasure soundness from inflatability]
\label{thm:erasure-from-inflatable}
Let $(\Phi, \gamma)$ be $\MM$-sound and suppose $R = \Phi \cup
\{\mathrm{nnf}(\neg\gamma)\}$ is $\Sigma_0$-inflatable.  If
$\erase{\Sigma_0}(\Phi) \folproves \erase{\Sigma_0}(\gamma)$, then
the pending goal is valid in $\MM$.
\end{theorem}

\begin{proof}
Suppose the goal is not valid.  By $\MM$-soundness there is an
expansion $\FM$ of $\FMstar$ with $\FM \folmodels \Phi$, and by the
contrapositive of \cref{def:si}(ii), $\FM \not\folmodels \gamma$;
hence $R$ is satisfiable.  By inflatability, $R$ has a model with
all $\Sigma_0$-guards total, which by \cref{lem:erasure-total} is a
model of $\erase{\Sigma_0}(R) = \erase{\Sigma_0}(\Phi) \cup
\{\mathrm{nnf}(\neg\erase{\Sigma_0}(\gamma))\}$ --- contradicting
the assumed provability.
\end{proof}

The rest of the section establishes checkable sufficient conditions
for inflatability, by passing through many-sorted logic.

\subsection{Sorted cores and relativisation}

\begin{definition}[Sorted core]
\label{def:sorted-core}
A totally specialised problem $R$ over guard family $\mathcal{G}$
(indexed by $\mathcal{T}$, now regarded as a set of \emph{sorts}) is
a \emph{sorted core} if it splits as $R = R_{\mathrm{log}} \uplus
R_{\mathrm{ty}}$ such that:
\begin{enumerate}[label=(SC\arabic*),leftmargin=3.2em]
\item every quantifier in $R_{\mathrm{log}}$ and in the guarded
  prefixes of $R_{\mathrm{ty}}$ is guarded: $\forall x.\,
  (\neg\guard{\tau}(x) \vee \cdots)$ or $\exists x.\,
  (\guard{\tau}(x) \wedge \cdots)$; the sort of the bound variable
  is the guard's index;
\item there is a \emph{sort assignment}: each function symbol
  $f$ carries a signature $(\tup{\sigma}; \varsigma)$ over
  $\mathcal{T}$, each predicate symbol other than the guards a
  signature $(\tup{\sigma})$, such that every occurrence of every
  symbol in $R$ is sort-correct (arguments of the declared sorts,
  computed from variable sorts upward), and every $\approx$-literal
  compares terms of equal sort;
\item $R_{\mathrm{ty}}$ consists of, for each function symbol $f$
  with signature $(\tup{\sigma}; \varsigma)$, the \emph{typing
  axiom} $\forall \tup{x}.\, \bigl(\bigvee_i
  \neg\guard{\sigma_i}(x_i)\bigr) \vee
  \guard{\varsigma}(f(\tup{x}))$ (for $0$-ary $f$: the ground
  \emph{typing fact} $\guard{\varsigma}(f)$); and guard atoms occur
  in $R_{\mathrm{log}}$ only in the guard positions of (SC1);
\item every sort $\tau \in \mathcal{T}$ is \emph{witnessed}: there
  is a ground term $w_\tau$ that is sort-correct of sort $\tau$.
\end{enumerate}
The \emph{sorted companion} $R^{\mathrm{ms}}$ is the many-sorted
problem over sorts $\mathcal{T}$ (with the signatures of (SC2))
whose formulas are those of $R_{\mathrm{log}}$ with each guarded
quantifier replaced by the corresponding sorted quantifier
($\forall x{:}\tau$, $\exists x{:}\tau$) and guard literals deleted;
$R_{\mathrm{ty}}$ has no counterpart (it is absorbed into the sorted
signature discipline).  Many-sorted semantics is standard: nonempty
domains $(D_\tau)_{\tau \in \mathcal{T}}$, sort-respecting
interpretations, equality per sort.
\end{definition}

\begin{lemma}[Pipeline problems are sorted cores]
\label{lem:pipeline-sorted}
The totally specialised translation of a drop-policy monomorphised
first-order-core problem (\cref{lem:spec-total}), restricted to
premises whose quantified types are witnessed, is a sorted core.
\end{lemma}

\begin{proof}
(SC1): all quantifiers of the translation are guarded ((F2), (F5),
and the guarded closures of every axiom scheme in
\cref{def:translation,def:transl-problem}), and (S2) turns each
guard into a $\guard{\tau}$-atom.  (SC2): sorts are read off the
$\CICm$ types: each flattened symbol $d^{(n)}$ or fused constant
descends from a constant or lifted constant with a closed $\CICm$
type, which after monomorphisation is a closed iterated product
over non-product data types; its signature is the list of argument
types and result type.  Sort-correctness of occurrences follows
from well-typedness of the $\CICm$ terms being translated
(\cref{lem:compositionality} respects types by \ref{M:typing}),
and $\approx$-literals compare translations of $\CICm$ terms of one
type ((F4), levelled equalities; \cref{lem:no-type-eq} excludes
type-level equations).  (SC3): the specialised $\mathrm{Flat}$
axioms and typing atoms have exactly the displayed shapes.  (SC4)
is the restriction in the statement.
\end{proof}

\begin{lemma}[Witness truth]
\label{lem:witness}
In every model $\FM$ of a sorted core $R$, and for every $\tau \in
\mathcal{T}$: $\guard{\tau}^{\FM} \neq \emptyset$; indeed the value
of any ground sort-correct term of sort $\tau$ lies in
$\guard{\tau}^{\FM}$.
\end{lemma}

\begin{proof}
By induction on the ground term, using the typing facts for
constants and the typing axioms for compound terms (their guard
premises hold by the induction hypothesis).
\end{proof}

\begin{theorem}[Relativisation]
\label{thm:relativisation}
A sorted core $R$ is satisfiable iff its sorted companion
$R^{\mathrm{ms}}$ is satisfiable.  Moreover, from a model of $R$
one obtains a model of $R^{\mathrm{ms}}$ with domains $D_\tau =
\guard{\tau}^{\FM}$, and conversely from a model of
$R^{\mathrm{ms}}$ a model of $R$ whose guard extensions are
(disjoint copies of) the sorted domains.
\end{theorem}

\begin{proof}
($\Rightarrow$)  Let $\FM \folmodels R$ with domain $\Dom$.  Set
$D_\tau \coloneqq \guard{\tau}^{\FM}$, nonempty by
\cref{lem:witness}.  For $f$ with signature $(\tup{\sigma};
\varsigma)$ define $f^{\mathrm{ms}}$ as the restriction of
$f^{\FM}$ to $\prod_i D_{\sigma_i}$; the typing axiom of $f$ (true
in $\FM$) guarantees values in $D_\varsigma$, so this is a legal
sorted interpretation.  Restrict predicates likewise.  We claim: for
every subformula $\varphi$ of a formula of $R_{\mathrm{log}}$ and
every sorted valuation $\xi$ (which is also an $\FM$-valuation),
$\FM, \xi \folmodels \varphi$ iff $\mathcal{N}, \xi \folmodels
\varphi^{\mathrm{ms}}$, where $\mathcal{N}$ is the sorted structure
just defined.  Induction on $\varphi$: atoms and their negations are
interpreted by restrictions, and all argument values lie in the
right domains by (SC2) and an inner induction on terms; the
connective cases are trivial; a guarded $\forall x.\,
(\neg\guard{\tau}(x) \vee \psi)$ holds in $\FM$ iff $\psi$ holds for
all $x$-values in $\guard{\tau}^{\FM} = D_\tau$, which by the
induction hypothesis is the truth condition of $\forall x{:}\tau.\,
\psi^{\mathrm{ms}}$; dually for $\exists$.  Hence $\mathcal{N}
\folmodels R^{\mathrm{ms}}$.

($\Leftarrow$)  Let $\mathcal{N} \folmodels R^{\mathrm{ms}}$ with
domains $(D_\tau)$.  Let $\Dom \coloneqq \biguplus_\tau D_\tau$
(tagged disjoint union; write $\iota_\tau : D_\tau \to \Dom$ for
the injections).  Interpret $\guard{\tau}^{\FM} \coloneqq
\iota_\tau(D_\tau)$; for $f$ with signature $(\tup{\sigma};
\varsigma)$ let $f^{\FM}(\tup{u}) \coloneqq
\iota_\varsigma(f^{\mathcal{N}}(\tup{d}))$ when $u_i =
\iota_{\sigma_i}(d_i)$ for all $i$, and an arbitrary fixed element
of $\Dom$ otherwise; predicates: true on sort-matching tuples
exactly when true in $\mathcal{N}$, false otherwise.  The typing
axioms hold by construction (guard-passing arguments are
sort-matching).  For $R_{\mathrm{log}}$, the same induction as
above, run in reverse, shows each formula true in $\FM$: guarded
quantifiers range over $\iota_\tau(D_\tau)$, which is a bijective
copy of $D_\tau$; sort-correct terms under sort-conforming
valuations never take the default branch; $\approx$-literals
compare values within one $\iota_\tau$-image, where $\iota_\tau$ is
injective.
\end{proof}

We also record the standard corollary used below.

\begin{corollary}[Countable models, many-sorted]
\label{cor:ls}
If a sorted problem $T$ over a countable signature and finitely
many sorts (all interpreted nonempty) is satisfiable, it has a
model with all domains countable.
\end{corollary}

\begin{proof}
Encode $T$ as a guarded unsorted problem by reversing the
construction of \cref{thm:relativisation} (guard each quantifier,
add typing axioms and one witness constant per sort --- this is the
classical relativisation, and the two directions proved above did
not use any property of $R$ beyond the sorted-core shape, so they
apply).  The unsorted problem is satisfiable, hence by the
L\"owenheim--Skolem theorem for countable signatures
\cite{ChangKeisler1990} it has a countable model; transferring back
through
\cref{thm:relativisation}($\Rightarrow$) yields a sorted model
whose domains are subsets of a countable set, and they are nonempty
by \cref{lem:witness}.
\end{proof}

\subsection{Monotonicity}

\begin{definition}[Monotonic sort {\cite[Def.~5.1]{Blanchette2016}}]
\label{def:monotonic}
Let $T$ be a many-sorted problem and $S$ a set of its sorts.  $S$
is \emph{monotonic} in $T$ if for every model $\mathcal{N}$ of $T$
there is a model $\mathcal{N}'$ of $T$ such that $D'_\tau$ is
infinite for every $\tau \in S$ and $D'_\rho = D_\rho$ for $\rho
\notin S$.
\end{definition}

\begin{definition}[Naked variables {\cite[Def.~5.6]{Blanchette2016}}]
\label{def:naked}
For NNF formulas: $\mathrm{NV}(p(\tup{t})) =
\mathrm{NV}(\neg p(\tup{t})) = \mathrm{NV}(t_1 \not\approx t_2) =
\emptyset$; $\mathrm{NV}(t_1 \approx t_2) = \{t_1, t_2\} \cap
\mathcal{V}$ (the variables among the two sides);
$\mathrm{NV}(\varphi \wedge \psi) = \mathrm{NV}(\varphi \vee \psi)
= \mathrm{NV}(\varphi) \cup \mathrm{NV}(\psi)$;
$\mathrm{NV}(\forall x{:}\tau.\, \varphi) = \mathrm{NV}(\varphi)$;
$\mathrm{NV}(\exists x{:}\tau.\, \varphi) = \mathrm{NV}(\varphi)
\setminus \{x\}$.  For a problem, $\mathrm{NV}(T) =
\bigcup_{\varphi \in T} \mathrm{NV}(\varphi)$.  A sort $\tau$ is
\emph{NV-safe} in $T$ if no variable of sort $\tau$ belongs to
$\mathrm{NV}(T)$.
\end{definition}

\begin{lemma}[Cloning; adapted from {\cite[Thm.~5.9]{Blanchette2016}} and {\cite[\S2.3]{Claessen2011}}]
\label{lem:cloning}
Let $\tau$ be NV-safe in the sorted problem $T$ and let
$\mathcal{N} \folmodels T$.  For any infinite cardinal $\kappa \geq
|D_\tau|$ there is $\mathcal{N}' \folmodels T$ with $|D'_\tau| =
\kappa$ and $D'_\rho = D_\rho$ for all $\rho \neq \tau$.
\end{lemma}

\begin{proof}
Let $E$ be a set of fresh elements with $|D_\tau \uplus E| =
\kappa$ and put $D'_\tau \coloneqq D_\tau \uplus E$, $D'_\rho
\coloneqq D_\rho$ otherwise.  Fix $e_0 \in D_\tau$ (nonempty) and
define the retraction $r : D'_\tau \to D_\tau$ by $r(e) = e_0$ for
$e \in E$ and $r(d) = d$ for $d \in D_\tau$; extend $r$ to other
sorts as the identity.  Interpret every function and predicate
symbol by precomposition with $r$:
$f^{\mathcal{N}'}(\tup{d}) \coloneqq f^{\mathcal{N}}(r(\tup{d}))$,
$p^{\mathcal{N}'}(\tup{d}) :\iff p^{\mathcal{N}}(r(\tup{d}))$.
For a valuation $\xi'$ over $\mathcal{N}'$ let $\xi \coloneqq r
\circ \xi'$.

\emph{Fact 1.}  For every term $t$ that is not a variable of sort
$\tau$: $\mathrm{val}_{\mathcal{N}',\xi'}(t) =
\mathrm{val}_{\mathcal{N},\xi}(t)$; and for a variable $x$ of sort
$\tau$: $r(\mathrm{val}_{\mathcal{N}',\xi'}(x)) =
\mathrm{val}_{\mathcal{N},\xi}(x)$.  (Induction on $t$; function
values are computed through $r$, so they land in the original
domains and agree.)

\emph{Fact 2.}  For every NNF formula $\varphi$ and every $\xi'$
such that $\xi'(x) \in D_\tau$ (no clone) for each free $x \in
\mathrm{NV}(\varphi)$ of sort $\tau$:\ if $\mathcal{N}, \xi
\folmodels \varphi$ then $\mathcal{N}', \xi' \folmodels \varphi$.
Induction on $\varphi$:
\begin{itemize}
\item Non-equality literals: all argument values retract correctly
  (Fact~1 and $r \circ \xi' = \xi$), and the interpretations factor
  through $r$, so the truth values agree exactly.
\item $t_1 \approx t_2$ (positive): each side is either a non-variable
  or a variable of another sort (then Fact~1 gives equality of
  values with the $\mathcal{N}$-side) or a naked variable of sort
  $\tau$ --- which by the valuation condition is mapped to a
  non-clone, so again its value coincides with its
  $\mathcal{N}$-value.  Hence the equation transfers.
\item $t_1 \not\approx t_2$: let $v'_i$ be the
  $\mathcal{N}'$-values and $v_i$ the $\mathcal{N}$-values; by
  Fact~1, $r(v'_i) = v_i$ (for non-variables and other-sort
  variables even $v'_i = v_i$).  From $v_1 \neq v_2$ and $r(v'_1)
  = v_1 \neq v_2 = r(v'_2)$ we get $v'_1 \neq v'_2$ since $r$ is a
  function.
\item $\wedge, \vee$: immediate.
\item $\forall x{:}\sigma.\, \psi$: for arbitrary $d \in
  D'_\sigma$, the retracted valuation is $\xi[x \mapsto r(d)]$ with
  $r(d) \in D_\sigma$, so $\mathcal{N}, \xi[x \mapsto r(d)]
  \folmodels \psi$; the valuation condition for $\psi$ under
  $\xi'[x \mapsto d]$ holds: if $x \in \mathrm{NV}(\psi)$ had sort
  $\tau$ then $x \in \mathrm{NV}(\forall x.\psi) =
  \mathrm{NV}(\psi)$ would make $\tau$ NV-unsafe in $T$,
  contradiction --- so either $x$'s sort is not $\tau$ or $x \notin
  \mathrm{NV}(\psi)$, and other free naked variables keep their
  values.  Apply the induction hypothesis.
\item $\exists x{:}\sigma.\, \psi$: take the $\mathcal{N}$-witness
  $d \in D_\sigma \subseteq D'_\sigma$; then $\xi'[x \mapsto d]$
  maps $x$ to a non-clone, so the valuation condition holds for
  $\psi$ even if $x \in \mathrm{NV}(\psi)$, and $r(d) = d$ keeps
  the retracted valuation aligned.  Apply the induction hypothesis.
\end{itemize}
Since sentences have no free variables, Fact~2 gives $\mathcal{N}'
\folmodels T$.
\end{proof}

\begin{lemma}[Forced infinity is monotonic]
\label{lem:forced-infinite}
If every model of $T$ interprets $\tau$ by an infinite domain, then
$\{\tau\}$ is monotonic in $T$ (take $\mathcal{N}' = \mathcal{N}$).
\qed
\end{lemma}

\begin{lemma}[Inductive infinity]
\label{lem:ind-infinity}
Let $R$ be a sorted core arising from the pipeline
(\cref{lem:pipeline-sorted}) and let $\tau$ be the sort of a
monomorphic instance $I\,\tup{\sigma}$ of a data inductive type.
Suppose $R$ contains, for some constructor $c_j$ of $I$ with a
recursive argument position $l_0$:
\begin{enumerate}[label=(\roman*)]
\item the specialised instance at $\tup{\sigma}$ of $c_j$'s
  injectivity axiom;
\item witnesses (SC4) for the sorts of all non-recursive argument
  positions of $c_j$ at $\tup{\sigma}$, and for $\tau$ itself a
  witness $w_\tau$ whose head constructor is some $c_i$ with $i
  \neq j$, together with the specialised instance at $\tup{\sigma}$
  of the discrimination axiom for the pair $(c_i, c_j)$.
\end{enumerate}
Then $D_\tau$ is infinite in every model of the sorted companion
$R^{\mathrm{ms}}$.
\end{lemma}

\begin{proof}
Fix a model; let $\tup{w}$ be the values of the non-recursive
witnesses and define $h : D_\tau \to D_\tau$ by $h(a) \coloneqq
c_j^{\mathrm{ms}}(\tup{w}, a)$ ($a$ in position $l_0$).  By the
injectivity axiom (instantiated at the fixed $\tup{w}$), $h$ is
injective.  By the discrimination axiom, the value of $w_\tau$
(a $c_i$-headed ground term) is not in the image of
$c_j^{\mathrm{ms}}$, in particular not in the image of $h$.  A set
admitting an injective, non-surjective self-map is
Dedekind-infinite.
\end{proof}

Whether the hypotheses of \cref{lem:ind-infinity} hold is a simple
syntactic check on the environment and the instance set: $I$ has a
constructor with a recursive argument whose non-recursive argument
types (at $\tup{\sigma}$) have closed inhabitants --- e.g.\
$\mathsf{nat}$; $\mathsf{list}\,\sigma$ for witnessed $\sigma$;
trees --- while $\mathsf{unit}$, $\mathsf{bool}$, enumerations, and
$\mathsf{list}$ of an empty type all fail it, as they must.  (If
every closed inhabitant of $\tau$ had head $c_j$, an innermost such
inhabitant would need a smaller closed inhabitant at the recursive
position, so a witnessed type with a recursive constructor always
has a second, non-recursive constructor to discriminate against;
condition (ii) is thus no real restriction beyond witnessing.)
Note that the criterion is \emph{internal}: it demands that the
injectivity and discrimination axioms be \emph{present in the
problem} --- which the translation guarantees for inductive types
occurring in it --- and therefore avoids the soundness gap of
infinity-by-introspection reported for Sledgehammer, where a type
known infinite in the background theory but not in the problem
makes the encoding unsound relative to the problem alone
\cite[\S7.4]{Blanchette2016}.

\subsection{The erasure theorem}

\begin{definition}[Erasable types]
\label{def:erasable}
Let $R$ be a sorted core.  A sort $\tau$ is \emph{erasable} if it
is witnessed (SC4) and either (a) NV-safe in $R^{\mathrm{ms}}$
(\cref{def:naked}), or (b) forced infinite in every model of
$R^{\mathrm{ms}}$ by \cref{lem:ind-infinity} (or any other
argument).  Write $\Sigma_0$ for any chosen set of erasable sorts.
\end{definition}

\begin{theorem}[Guard-erasure soundness]
\label{thm:erasure-main}
Let $R = \Phi \cup \{\mathrm{nnf}(\neg\gamma)\}$ be a sorted core
and $\Sigma_0$ a set of erasable sorts.  Then $R$ is
$\Sigma_0$-inflatable.  Consequently
(\cref{thm:erasure-from-inflatable}), if $(\Phi, \gamma)$ is
$\MM$-sound and $\erase{\Sigma_0}(\Phi) \folproves
\erase{\Sigma_0}(\gamma)$, the goal is valid in $\MM$; and by
\cref{prop:erasure-complete} no first-order proof is lost.
\end{theorem}

\begin{proof}
Assume $R$ satisfiable; we construct a model of $R$ with all
$\Sigma_0$-guards total.

\emph{Step 1 (sorted, countable).}  By
\cref{thm:relativisation}($\Rightarrow$), $R^{\mathrm{ms}}$ is
satisfiable; by \cref{cor:ls} it has a model $\mathcal{N}_1$ with
all domains countable (nonempty).

\emph{Step 2 (inflate the erasable sorts).}  Enumerate $\Sigma_0 =
\{\tau_1, \ldots, \tau_m\}$ (finite).  Define $\mathcal{N}_2$ by
treating each $\tau_i$ in turn: if $\tau_i$ is forced infinite, its
domain is already countably infinite and we do nothing; otherwise
$\tau_i$ is NV-safe, and we apply \cref{lem:cloning} with $\kappa =
\aleph_0$, which changes only $D_{\tau_i}$.  NV-safety is a
property of the problem $R^{\mathrm{ms}}$, not of the model, so
successive applications remain justified; forced-infinite sorts
keep their (unchanged, infinite) domains.  After the $m$ steps,
every $\tau \in \Sigma_0$ has $|D_\tau| = \aleph_0$ and all other
domains are unchanged, countable and nonempty.

\emph{Step 3 (merge).}  Let $U$ be a fixed countably infinite set.
Choose bijections $b_\tau : D_\tau \to U$ for $\tau \in \Sigma_0$
and injections $j_\rho : D_\rho \hookrightarrow U$ for $\rho \in
\mathcal{T} \setminus \Sigma_0$ (possible: the domains are
countable).  Write $\mathrm{emb}_\sigma$ for $b_\sigma$ or
$j_\sigma$ as appropriate, and $\mathrm{proj}_\sigma$ for the
(partial) inverse, total for $\sigma \in \Sigma_0$.  Define the
unsorted structure $\FM$ with domain $U$:
\begin{itemize}
\item $\guard{\rho}^{\FM} \coloneqq j_\rho(D_\rho)$ for $\rho
  \notin \Sigma_0$; \ $\guard{\tau}^{\FM} \coloneqq U$ for $\tau
  \in \Sigma_0$;
\item for $f$ with signature $(\tup{\sigma}; \varsigma)$:
  $f^{\FM}(\tup{u}) \coloneqq
  \mathrm{emb}_\varsigma\bigl(f^{\mathcal{N}_2}(
  \mathrm{proj}_{\sigma_1}(u_1), \ldots)\bigr)$ when every
  $\mathrm{proj}_{\sigma_i}(u_i)$ is defined, and
  $\mathrm{emb}_\varsigma(d^0_\varsigma)$ (a fixed default
  $d^0_\varsigma \in D_\varsigma$) otherwise;
\item for a predicate $p$ with signature $(\tup{\sigma})$:
  $p^{\FM}(\tup{u})$ iff every $\mathrm{proj}_{\sigma_i}(u_i)$ is
  defined and $p^{\mathcal{N}_2}$ holds of the projections.
\end{itemize}
Note $\mathrm{proj}_\sigma \circ \mathrm{emb}_\sigma =
\mathrm{id}$ for every sort, and $\mathrm{proj}_\tau$ is total and
bijective for $\tau \in \Sigma_0$ --- this totality is exactly what
Step~2 bought.

\emph{Claim: $\FM \folmodels R$, with all $\Sigma_0$-guards total.}
Totality of the $\Sigma_0$-guards holds by definition.  For the
formulas, we prove: for every formula $\varphi \in
R_{\mathrm{log}}$, every subformula $\psi$ of $\varphi$, and every
sort-conforming pair of valuations ($\xi$ over $\mathcal{N}_2$,
$\nu$ over $\FM$) with $\nu(x) =
\mathrm{emb}_{\mathrm{sort}(x)}(\xi(x))$ for the free variables:
\[
  \mathcal{N}_2, \xi \folmodels \psi^{\mathrm{ms}}
  \quad\Longleftrightarrow\quad
  \FM, \nu \folmodels \psi .
\]
Terms first: by induction, the $\FM$-value of a sort-$\varsigma$
term $t$ under $\nu$ is $\mathrm{emb}_\varsigma$ of its
$\mathcal{N}_2$-value under $\xi$ (variables by conformity;
applications because $\mathrm{proj} \circ \mathrm{emb} =
\mathrm{id}$ keeps every projection defined, so the default branch
is never taken).  Atoms: $p(\tup{t})$ transfers by the same
identity; $t_1 \approx t_2$ at sort $\varsigma$ transfers because
$\mathrm{emb}_\varsigma$ is injective.  Connectives: trivial.
Quantifier at a sort $\rho \notin \Sigma_0$ (guard kept): the
$\FM$-side quantifier is guarded by $\guard{\rho}^{\FM} =
j_\rho(D_\rho)$, a bijective copy of the $\mathcal{N}_2$-range ---
both sides range over matching values and the induction hypothesis
applies.  Quantifier at a sort $\tau \in \Sigma_0$ (guard
\emph{present} in $R$; recall we are proving $\FM \folmodels R$,
not its erasure): the guard is total, so the $\FM$-side ranges over
all of $U = b_\tau(D_\tau)$ --- again a bijective copy of the
sorted range, and the induction hypothesis applies.  This proves
every formula of $R_{\mathrm{log}}$ true in $\FM$ (its
$^{\mathrm{ms}}$-image being true in $\mathcal{N}_2$).

$R_{\mathrm{ty}}$: for a typing axiom of $f : (\tup{\sigma};
\varsigma)$, let $\tup{u}$ pass the (remaining) guards.  For
$\sigma_i \in \Sigma_0$ the projection is total; for $\sigma_i
\notin \Sigma_0$ the guard $\guard{\sigma_i}(u_i)$ asserts $u_i \in
j_{\sigma_i}(D_{\sigma_i})$, so the projection is defined.  Hence
$f^{\FM}(\tup{u})$ is $\mathrm{emb}_\varsigma$ of a value in
$D_\varsigma$, and the conclusion $\guard{\varsigma}(f(\tup{u}))$
holds ($= U$ if $\varsigma \in \Sigma_0$, $=
j_\varsigma(D_\varsigma) \ni$ the value otherwise).  Typing facts
and witness terms: likewise.

So $\FM \folmodels R$ with total $\Sigma_0$-guards: $R$ is
$\Sigma_0$-inflatable.
\end{proof}

\subsection{Necessity of the side conditions}

\begin{proposition}[Witnessing is necessary]
\label{prop:nec-inhab}
There is an $\MM$-sound sorted-core problem $(\Phi, \gamma)$ with
$\gamma$'s pending goal \emph{invalid} in $\MM$, and a
\emph{non-witnessed} but NV-safe type $\tau$, such that
$\erase{\{\tau\}}(\Phi) \folproves \erase{\{\tau\}}(\gamma)$.
\end{proposition}

\begin{proof}
Let $\mathsf{empty}$ be the data inductive type with no
constructors, and take the single premise $\varphi = \Pi
x{:}\mathsf{empty}.\ \mathsf{O} \ceq_{\mathsf{nat}}
\mathsf{S}\,\mathsf{O}$ --- a \emph{valid} proposition
($\sem{\mathsf{empty}} = \emptyset$ by \ref{M:ind}, so the product
is vacuously $\holds$; it is also provable by
$\mathsf{empty}$-elimination in CIC).  Its specialised translation
is $\forall x.\ \neg\guard{\mathsf{empty}}(x) \vee
(\hat{\mathsf{O}} \approx \hat{\mathsf{S}}(\hat{\mathsf{O}}))$.
The type $\mathsf{empty}$ is NV-safe ($x$ occurs in no equation;
the inversion axiom for a constructorless type is $\forall x.\
\neg\guard{\mathsf{empty}}(x) \vee \bot$, equation-free), but not
witnessed.  Take the goal $\gamma$ translating $\varphi_G =
(\mathsf{O} \ceq \mathsf{S}\,\mathsf{O})$, invalid in $\MM$ by
\ref{M:ind}(c).  Erasing the $\mathsf{empty}$-guard turns the
premise into $\forall x.\ \hat{\mathsf{O}} \approx
\hat{\mathsf{S}}(\hat{\mathsf{O}})$, from which
$\erase{}(\gamma) = \gamma$ follows immediately.
\end{proof}

\begin{proposition}[Monotonicity is necessary]
\label{prop:nec-mono}
There is an $\MM$-sound sorted-core problem with an invalid pending
goal and a \emph{witnessed} but non-monotonic type $\tau$ such that
erasing the $\tau$-guards proves the goal.
\end{proposition}

\begin{proof}
Take $\tau = \mathsf{unit}$ (single constructor $\mathsf{tt}$) and
as premise its own inversion axiom, emitted by the translation
itself: $\forall x.\ \neg\guard{\mathsf{unit}}(x) \vee x \approx
\hat{\mathsf{tt}}$.  It is valid ($\sem{\mathsf{unit}} =
\{\sem{\mathsf{tt}}\}$).  The sort is witnessed
($\hat{\mathsf{tt}}$) but the variable $x$ occurs naked in a
positive equation, and indeed $\mathsf{unit}$ is genuinely
non-monotonic: the axiom bounds the guard extension by one element.
Erase the guard: $\forall x.\ x \approx \hat{\mathsf{tt}}$.
Instantiating $x$ by $\hat{\mathsf{O}}$ and by
$\hat{\mathsf{S}}(\hat{\mathsf{O}})$ yields $\hat{\mathsf{O}}
\approx \hat{\mathsf{S}}(\hat{\mathsf{O}})$: the invalid goal of
\cref{prop:nec-inhab} again follows.
\end{proof}

\begin{remark}[The case-scrutinee counterexample, revisited]
\label{rem:jar-counterexample}
\cite[\S5.6]{CzajkaKaliszyk2018} observes that the guards on the
free variables of a case scrutinee can never be omitted, exhibiting
a one-constructor type $c$ with $\mathsf{I}^0(c : \mathsf{Set}
\coloneqq c_0 : c)$ where omitting the guard produces the axiom
$\forall x.\ x = c_0 \wedge F\,x = c_0$ and thence an inconsistency.
Through the lens of this section the diagnosis is exact:
the conjunct $\forall x.\ x \approx c_0$ is the inversion content
of a \emph{singleton} --- witnessed, maximally non-monotonic ---
type, i.e.\ precisely the situation of \cref{prop:nec-mono}.  The
guards that \emph{may} be dropped are those of
\cref{def:erasable}; the case-scrutinee guards at singleton and
enumeration types are the paradigmatic non-examples, and
\cref{thm:erasure-main} never licenses them.  Conversely, the
folklore practice of omitting guards at types like $\mathsf{nat}$
or $\mathsf{list}\,\mathsf{nat}$ is licensed, for the first time
with a proof, by \cref{lem:ind-infinity} together with
\cref{thm:erasure-main}.
\end{remark}

\begin{remark}[Scope, and comparison with the featherweight encodings]
\label{rem:erasure-scope}
(i) Our erasure is \emph{per-type}, like the lightweight encodings
$\widetilde{\mathsf{g}}?$ of \cite{Claessen2011,Blanchette2016};
the featherweight $\widetilde{\mathsf{g}}??$ additionally drops
guards of non-naked \emph{variables} at non-monotonic types.  The
analogous per-variable refinement is meaningful in our setting, but
its correctness proof requires interpreting a kept guard as
\emph{false} on merged junk (the device of
\cite[Thm.~5.20]{Blanchette2016}) inside our single-domain merge,
where the same junk element is shared by all erased sorts; we leave
this refinement to future work.  (ii) The freedom to choose
$\Sigma_0$ as \emph{any} subset of the erasable sorts (cf.\
\cite[Rem.~5.21]{Blanchette2016}) is built into
\cref{thm:erasure-main}.  (iii) The theorem is proved for the
drop-policy monomorphic pipeline.  With kept generic premises the
problem is not a sorted core --- binary $\hastype$-atoms with
variable type arguments remain --- and erasure of the specialised
part is not covered; the appropriate tool there would be a
polymorphic monotonicity calculus in the style of
\cite[\S6]{Blanchette2016}, which we also leave open.
(iv) Erasure interacts with \cref{rem:funext}: after erasure at an
\emph{erasable} type, the collapse phenomena that make guard
omission unsound in the implemented translation cannot occur, since
erasable types admit models where the guard is total.
\end{remark}
